\documentclass[11pt,letterpaper]{article}

% A conservative, arXiv-friendly preamble. pdfLaTeX is the default; the same
% source also compiles with XeLaTeX without requiring external font files.
% Source references: https://arxiv.org/src/2501.18774,
% https://arxiv.org/src/2408.11653, and https://arxiv.org/src/1807.03761.
\usepackage{iftex}
\ifPDFTeX%
  \usepackage[utf8]{inputenc}
\fi%
\usepackage[T1]{fontenc}
\usepackage{lmodern}

\usepackage[margin=1in]{geometry}
\usepackage{microtype}
\usepackage{amsmath,amssymb,amsfonts,amsthm,mathtools,mathrsfs}
\usepackage{xcolor}
\usepackage{enumitem}
\usepackage{booktabs}
\usepackage{graphicx}
\usepackage{url}

% Keep links legible on screen and quiet in print.
\definecolor{linkblue}{RGB}{20,72,118}
\definecolor{citeburgundy}{RGB}{128,36,52}
\usepackage[
  colorlinks=true,
  linkcolor=linkblue,
  citecolor=citeburgundy,
  urlcolor=linkblue,
  bookmarks=true,
  bookmarksnumbered=true
]{hyperref}

% Alpoge-style mathematical structure: a single theorem counter within each
% section, with equations numbered by section as well.
\numberwithin{equation}{section}

\theoremstyle{plain}
\newtheorem{theorem}{Theorem}[section]
\newtheorem{proposition}[theorem]{Proposition}
\newtheorem{lemma}[theorem]{Lemma}
\newtheorem{corollary}[theorem]{Corollary}
\newtheorem{claim}[theorem]{Claim}
\newtheorem{conjecture}[theorem]{Conjecture}

\theoremstyle{definition}
\newtheorem{definition}[theorem]{Definition}
\newtheorem{example}[theorem]{Example}
\newtheorem{question}[theorem]{Question}
\newtheorem{assumption}[theorem]{Assumption}

\theoremstyle{remark}
\newtheorem{remark}[theorem]{Remark}

% Common operators and notation. Add paper-specific commands here rather than
% scattering definitions through the body.
\newcommand{\N}{\mathbb{N}}
\newcommand{\Z}{\mathbb{Z}}
\newcommand{\Q}{\mathbb{Q}}
\newcommand{\R}{\mathbb{R}}
\newcommand{\C}{\mathbb{C}}
\newcommand{\cO}{\mathcal{O}}
\DeclareMathOperator{\Hom}{Hom}
\DeclareMathOperator{\End}{End}
\DeclareMathOperator{\Gal}{Gal}
\DeclareMathOperator{\Spec}{Spec}

\setlist{itemsep=0.25em,topsep=0.5em}
\urlstyle{same}
\allowdisplaybreaks[2]
\emergencystretch=2em

\title{A Local-to-Global Principle for Structured Families}
\author{Author Name}
\date{}

\begin{document}

\maketitle

\begin{abstract}
We study a family of arithmetic objects equipped with compatible local
conditions. Under a finite-obstruction hypothesis, we show that local
solvability forces a global solution and obtain a quantitative refinement.
The proof separates the geometric input from the counting argument, making
the source of each hypothesis explicit. This abstract is starter text: replace
it with the question, principal result, method, and significance of the paper.
\end{abstract}

\section{Introduction}\label{sec:introduction}

Many existence problems in arithmetic begin with a comparison between global
objects and their local avatars. The central question is whether compatible
local data can be assembled into a global solution. In the setting considered
here, the obstruction is encoded by a finite set $\Omega$ and a map
\begin{equation}\label{eq:obstruction-map}
  \operatorname{ob}\colon \mathcal{X}\longrightarrow \prod_{v\in\Omega}
  \mathcal{X}_v.
\end{equation}
Our goal is to isolate hypotheses under which the image of~\eqref{eq:obstruction-map}
is large enough to detect global solvability.

The main result, Theorem~\ref{thm:main}, has two useful features. First, its
hypotheses are entirely local away from the finite set $\Omega$. Second, the
argument gives a lower bound rather than existence alone. This separation is
useful when the same geometric input is paired with different counting
methods.

\subsection{Outline of the paper}

Section~\ref{sec:framework} fixes notation and records the local hypotheses.
Section~\ref{sec:main-result} states the main theorem and organizes its proof
into two lemmas. Section~\ref{sec:questions} collects variants and open
questions.

\section{The framework}\label{sec:framework}

Let $K$ be a field, let $V$ be a geometrically integral $K$-variety, and let
$\mathcal{H}$ be a countable family of height-bounded subsets of $V(K)$. Write
$\mathcal{H}(B)$ for the members of height at most $B$.

\begin{definition}\label{def:admissible}
A datum $x\in\mathcal{H}(B)$ is \emph{admissible} if its localization $x_v$
lies in a prescribed open subset $U_v\subseteq V(K_v)$ for every
$v\in\Omega$.
\end{definition}

We use the following model hypothesis to distinguish the structural input from
the analytic one.

\begin{assumption}\label{ass:equidistribution}
There are constants $\delta>0$ and $\eta>0$ such that, uniformly over
admissible residue conditions,
\begin{equation}\label{eq:counting}
  \#\mathcal{H}(B)=\delta B+O\!\left(B^{1-\eta}\right).
\end{equation}
\end{assumption}

\begin{remark}
In an application, $\delta$ is normally a product of local densities. The
substantive task is to prove that the product is nonzero and that the error
in~\eqref{eq:counting} is uniform.
\end{remark}

\section{The main result}\label{sec:main-result}

\begin{theorem}\label{thm:main}
Suppose that Assumption~\ref{ass:equidistribution} holds and that the local
conditions $U_v$ have positive density. Then there is a constant $c>0$ such
that, for all sufficiently large $B$,
\begin{equation}\label{eq:main-bound}
  \#\{x\in\mathcal{H}(B):x\text{ is admissible}\}\geq cB.
\end{equation}
In particular, infinitely many admissible data exist.
\end{theorem}

The proof is deliberately short in this template; a real paper should state
where each uniformity estimate enters.

\begin{lemma}\label{lem:local-density}
The combined local condition has density
$\delta_{\Omega}=\prod_{v\in\Omega}\delta_v>0$.
\end{lemma}

\begin{proof}
The set $\Omega$ is finite and every factor $\delta_v$ is positive, so their
product is positive.
\end{proof}

\begin{proof}[Proof of Theorem~\ref{thm:main}]
Apply Assumption~\ref{ass:equidistribution} after imposing the residue
conditions defining admissibility. Lemma~\ref{lem:local-density} gives
\begin{equation*}
  \#\{x\in\mathcal{H}(B):x\text{ is admissible}\}
  =\delta_{\Omega}B+O\!\left(B^{1-\eta}\right).
\end{equation*}
For sufficiently large $B$, the error is at most
$\tfrac12\delta_{\Omega}B$. Thus~\eqref{eq:main-bound} holds with
$c=\tfrac12\delta_{\Omega}$.
\end{proof}

\begin{corollary}\label{cor:existence}
Under the hypotheses of Theorem~\ref{thm:main}, at least one admissible global
datum exists.
\end{corollary}

\begin{proof}
Take $B$ large enough that the right-hand side
of~\eqref{eq:main-bound} is positive.
\end{proof}

\section{Further questions}\label{sec:questions}

The argument suggests three natural directions.

\begin{enumerate}
  \item Can the finite set $\Omega$ vary with the height bound $B$?
  \item Can the power-saving error in~\eqref{eq:counting} be made effective?
  \item Which geometric hypotheses guarantee positive local density?
\end{enumerate}

\begin{question}
Is there an intrinsic description of the smallest obstruction set $\Omega$
for which Theorem~\ref{thm:main} remains valid?
\end{question}

For a model of concise arithmetic exposition and standard background on local
fields, see Serre~\cite{serre-local-fields}. Replace this citation with the
sources actually used in the paper.


\subsection*{Acknowledgments}
Add acknowledgments, grant support, and data or code availability statements
here.

\bibliographystyle{amsalpha}
\bibliography{references}

\end{document}
